%%======================================================================
%%  NT-2: Literature Extraction & Convention Fixing
%%        for Entire-Function Property and Ghost-Freedom Criteria
%%
%%  This document collects formulas from the literature needed to
%%  establish the entire-function property of the SCT form factors
%%  F_1(z), F_2(z) and collect the literature ghost-freedom criteria
%%  against which the physical propagator must be tested.
%%
%%  RESEARCH DOCUMENT: formulas extracted from cited papers.
%%  No original derivations --- gaps noted explicitly.
%%======================================================================
\documentclass[12pt,a4paper]{article}

%% ---- packages -------------------------------------------------------
\usepackage[utf8]{inputenc}
\usepackage[T1]{fontenc}
\usepackage{lmodern}
\usepackage{amsmath,amssymb,amsthm,mathtools}
\usepackage[margin=2.5cm]{geometry}
\usepackage{booktabs}
\usepackage{enumitem}
\usepackage{hyperref}
\usepackage{xcolor}
\usepackage{fancyhdr}
\usepackage{titlesec}
\usepackage{longtable}

%% ---- theorem-like environments --------------------------------------
\newtheorem{proposition}{Proposition}[section]
\newtheorem{lemma}[proposition]{Lemma}
\newtheorem{corollary}[proposition]{Corollary}
\newtheorem{theorem}[proposition]{Theorem}
\theoremstyle{definition}
\newtheorem{definition}[proposition]{Definition}
\newtheorem{convention}[proposition]{Convention}
\theoremstyle{remark}
\newtheorem{remark}[proposition]{Remark}

%% ---- custom commands ------------------------------------------------
\newcommand{\Tr}{\operatorname{Tr}}
\newcommand{\tr}{\operatorname{tr}}
\newcommand{\Rsq}{R_{\mu\nu\rho\sigma}R^{\mu\nu\rho\sigma}}
\newcommand{\Ricsq}{R_{\mu\nu}R^{\mu\nu}}
\newcommand{\Weylsq}{C_{\mu\nu\rho\sigma}C^{\mu\nu\rho\sigma}}
\newcommand{\dd}{\mathrm{d}}
\newcommand{\ii}{\mathrm{i}}
\newcommand{\Id}{\mathrm{Id}}
\newcommand{\erfi}{\operatorname{erfi}}
\DeclareMathOperator{\sgn}{sgn}
\DeclareMathOperator{\ord}{ord}

%% ---- annotation commands --------------------------------------------
\newcommand{\fromlit}[1]{\textcolor{blue!70!black}{\textsf{[#1]}}}
\newcommand{\oursign}{\textcolor{teal!80!black}{\textsf{[our convention]}}}
\newcommand{\gap}{\textcolor{red!80!black}{\textsf{[GAP]}}}
\newcommand{\needscheck}{\textcolor{orange!80!black}{\textsf{[needs check]}}}
\newcommand{\exactlit}{\textcolor{blue!60!black}{\textsf{[exact from lit.]}}}
\newcommand{\verified}[1]{%
  \par\smallskip\noindent
  \colorbox{green!10}{\parbox{\dimexpr\linewidth-2\fboxsep}{%
    \textcolor{green!50!black}{\textbf{[Verified]}} #1}}%
  \par\smallskip}
\newcommand{\signcorrection}[1]{%
  \par\noindent\colorbox{red!10}{%
    \parbox{\dimexpr\linewidth-2\fboxsep}{%
      \textcolor{red!50!black}{\textbf{[Sign correction]}}
      \textcolor{red!50!black}{#1}%
    }%
  }\par%
}

%% ---- page style -----------------------------------------------------
\pagestyle{fancy}
\fancyhf{}
\fancyhead[L]{\small\textsc{NT-2: Literature \& Conventions for Entire-Function Property}}
\fancyhead[R]{\small\thepage}
\renewcommand{\headrulewidth}{0.4pt}

%% ---- title ----------------------------------------------------------
\title{\textbf{NT-2: Literature Extraction \& Convention Fixing\\[4pt]
       for Entire-Function Property and Ghost-Freedom Criteria}}
\author{David Alfyorov}
\date{March 2026 --- Working Document}

%%======================================================================
\begin{document}
\maketitle

\begin{center}
\small\textit{Credits: Aliaksandr Samatyia contributed research-assistance and workflow support.}
\end{center}

\thispagestyle{empty}

\begin{abstract}
This document collects all formulas, definitions, and analytic results
from the literature needed to establish two logically distinct parts of
the NT-2 problem: (i)~the nonlocal form factors $F_1(z)$ and $F_2(z)$
should be analysed as \textbf{entire functions} of their argument
$z = \Box/\Lambda^2$; (ii)~the physical propagator denominators must be
checked against the standard \textbf{ghost-freedom criteria} used in the
nonlocal-gravity literature. All formulas are extracted from the
literature; no original derivations are performed. Each formula is
annotated with its source, equation number, and sign convention status.
Gaps are noted explicitly.

\medskip\noindent
\textbf{Important scope note.}
This literature document records the benchmark criteria for
ghost-freedom; it does \emph{not} by itself establish that the repaired
SCT propagator satisfies them. The current NT-4a normalization in fact
shows a positive-real-axis zero in the TT denominator, so the
ghost-freedom benchmark is presently \emph{not} met in that
normalization.

\medskip\noindent
\textbf{Key references:}
\begin{itemize}[nosep]
\item Codello--Zanusso (CZ): arXiv:1203.2034 [hep-th].
\item Codello--Percacci--Rahmede (CPR): arXiv:0805.2909 [hep-th].
\item Barvinsky--Vilkovisky (BV): Nucl.\ Phys.\ B \textbf{282} (1987) 163;
  B \textbf{333} (1990) 471, 512.
\item Tomboulis: hep-th/9702080.
\item Biswas--Mazumdar--Siegel (BMS): hep-th/0508194.
\item Modesto: arXiv:1107.2403 [hep-th].
\item Vassilevich: hep-th/0306138.
\item Avramidi: hep-th/9704166.
\item Hadamard: \emph{Th\`ese}, 1893 (entire function factorization).
\item Boas: \emph{Entire Functions}, Academic Press, 1954.
\end{itemize}
\end{abstract}

\tableofcontents
\newpage

%%======================================================================
\section{Background: From NT-1 to NT-2}
\label{sec:background}
%%======================================================================

The NT-1 derivation established the per-spin nonlocal form factors
$h_C^{(s)}(z)$ and $h_R^{(s)}(z)$ for the SCT spectral action, and the
NT-1b phases assembled the Standard Model totals:
\begin{equation}
  F_i(z) = \frac{1}{16\pi^2}\sum_{s=0,1/2,1} N_s\, h_i^{(s)}(z),
  \qquad i \in \{C, R\}.
  \label{eq:Fi-definition}
\end{equation}
All per-spin form factors are built from the \textbf{master function}
\begin{equation}
  \boxed{
  \phi(z) = \int_0^1 e^{-\xi(1-\xi)z}\,\dd\xi
  = e^{-z/4}\sqrt{\frac{\pi}{z}}\,\erfi\!\left(\frac{\sqrt{z}}{2}\right)
  }
  \label{eq:phi-master}
\end{equation}
together with rational functions of $z$ whose coefficients are
determined by spin and field content
\fromlit{CZ Eqs.~(16)--(21); NT-1 Theorem~4.1}.

\begin{remark}[Status from NT-1/NT-1b]
\label{rem:nt1-status}
The NT-1/NT-1b derivation chain delivered:
\begin{enumerate}[nosep]
  \item Closed-form expressions for all six form factors
    $h_C^{(0)}, h_R^{(0)}, h_C^{(1/2)}, h_R^{(1/2)},
     h_C^{(1)}, h_R^{(1)}$;
  \item Phase~3 SM totals $\alpha_C = 13/120$,
    $\alpha_R(\xi) = 2(\xi-1/6)^2$;
  \item Numerical verification at 100-digit precision (274+ checks).
\end{enumerate}
The NT-2 task is to prove these form factors are
\emph{entire functions of $z$}, establish their growth order/type,
and formulate the propagator benchmark that must be checked separately.
\par\smallskip
After the repaired NT-4a scan, this benchmark is known \emph{not} to be
satisfied by the current normalized TT denominator, which develops a
positive-real-axis zero near $z \approx 2.41484$.
\end{remark}


%%======================================================================
\section{The Master Function: Integral Representation}
\label{sec:master-function}
%%======================================================================

\subsection{Definitions from the literature}

\begin{definition}[Master function $\phi(z)$]
\label{def:phi}
\exactlit\;
The master function appearing in the CZ form factor calculus is
\fromlit{CZ Eq.~(14)}:
\begin{equation}
  \phi(z) = \int_0^1 e^{-\xi(1-\xi)z}\,\dd\xi\,.
  \label{eq:phi-integral}
\end{equation}
\end{definition}

\begin{remark}[Alternative representations]
Two equivalent closed forms are known \fromlit{CZ Eq.~(14) footnote;
BV 1990}:
\begin{align}
  \phi(z) &= e^{-z/4}\sqrt{\frac{\pi}{z}}\,\erfi\!\left(
    \frac{\sqrt{z}}{2}\right)
    \label{eq:phi-erfi}\\
  &= e^{-z/4}\sqrt{\frac{\pi}{z}}\,\frac{2}{\sqrt{\pi}}
    \int_0^{\sqrt{z}/2} e^{t^2}\,\dd t\,.
    \label{eq:phi-dawson}
\end{align}
The second form uses the definition $\erfi(w) =
(2/\sqrt{\pi})\int_0^w e^{t^2}\,\dd t$.
\end{remark}

\begin{proposition}[Taylor expansion of $\phi$]
\label{prop:phi-taylor}
\exactlit\;
From the integral representation \eqref{eq:phi-integral}:
\begin{equation}
  \phi(z) = \sum_{n=0}^{\infty} \frac{(-1)^n\, n!}{(2n+1)!}\, z^n\,.
  \label{eq:phi-taylor}
\end{equation}
\end{proposition}
\begin{proof}
Expand $e^{-\xi(1-\xi)z} = \sum_{n=0}^\infty \frac{(-z)^n}{n!}
[\xi(1-\xi)]^n$ and integrate term-by-term (justified by
uniform convergence on $[0,1]$ for any $z\in\mathbb{C}$):
\begin{equation}
  \int_0^1 [\xi(1-\xi)]^n\,\dd\xi = B(n+1,n+1)
  = \frac{(n!)^2}{(2n+1)!}\,,
\end{equation}
where $B$ is the Euler beta function.  Hence the $n$-th coefficient is
$(-1)^n (n!)^2 / [n!\,(2n+1)!] = (-1)^n n!/(2n+1)!$.
\end{proof}

\begin{corollary}[Key values]
\label{cor:phi-values}
\begin{align}
  \phi(0) &= 1, \label{eq:phi0}\\
  \phi'(0) &= -\frac{1}{6},\label{eq:phi-prime0}\\
  \phi''(0) &= \frac{1}{30}\label{eq:phi-double-prime0}\,.
\end{align}
These follow immediately from \eqref{eq:phi-taylor} with
$a_0 = 1$, $a_1 = -1/6$, $a_2 = 1/60$, so
$\phi''(0) = 2!\cdot a_2 = 1/30$.
\end{corollary}

\verified{$\phi(0)=1$, $\phi'(0)=-1/6$, $\phi''(0)=1/30$
confirmed by mpmath at 100-digit precision.  Identity
$\phi(0)=1$ also verified formally in Lean~4
(\texttt{PhiEntire.lean}, Theorem~\texttt{phi\_zero}).}


%%======================================================================
\section{Entire Functions: Definitions and Standard Results}
\label{sec:entire-theory}
%%======================================================================

\subsection{Basic definitions}

\begin{definition}[Entire function]
\label{def:entire}
\exactlit\;
A function $f:\mathbb{C}\to\mathbb{C}$ is \textbf{entire} if it is
holomorphic on all of $\mathbb{C}$, i.e., its Taylor series at the
origin converges everywhere:
\begin{equation}
  f(z) = \sum_{n=0}^\infty a_n z^n,
  \qquad \limsup_{n\to\infty} |a_n|^{1/n} = 0\,.
\end{equation}
\end{definition}

\begin{definition}[Order and type]
\label{def:order-type}
\exactlit\;
\fromlit{Boas, Ch.~2; Hadamard 1893}
For a non-constant entire function $f$, the \textbf{order} $\rho$ is
\begin{equation}
  \rho = \limsup_{r\to\infty}
  \frac{\log\log M(r)}{\log r}\,,
  \qquad M(r) = \max_{|z|=r}|f(z)|\,,
  \label{eq:order}
\end{equation}
and for $0 < \rho < \infty$, the \textbf{type} $\sigma$ is
\begin{equation}
  \sigma = \limsup_{r\to\infty}
  \frac{\log M(r)}{r^\rho}\,.
  \label{eq:type}
\end{equation}
An entire function of order~$\rho$ and type~$\sigma$ satisfies
$M(r) \leq C\, e^{(\sigma+\epsilon)r^\rho}$ for every $\epsilon > 0$.
\end{definition}

\begin{remark}[Terminology]
An entire function of order~1 is said to be of
\textbf{exponential type}.  For $\phi(z)$ we will establish
$\rho = 1$, $\sigma = 1/4$, so $\phi$ is entire of exponential
type~$1/4$.
\end{remark}


\subsection{Hadamard factorization theorem}

\begin{theorem}[Hadamard, 1893]
\label{thm:hadamard}
\exactlit\;
\fromlit{Hadamard 1893; Boas Ch.~2, Thm.~2.7.1}
Let $f$ be an entire function of finite order~$\rho$ with zeros
$\{z_n\}_{n=1}^\infty$ (counted with multiplicity), a zero of
order~$m$ at the origin, and genus
$p = \lfloor\rho\rfloor$.  Then
\begin{equation}
  \begin{aligned}
    f(z) &= z^m\, e^{P(z)} \\
         &\quad \times \prod_{n=1}^\infty E_p\!\left(\frac{z}{z_n}\right),
  \end{aligned}
  \label{eq:hadamard}
\end{equation}
where $P(z)$ is a polynomial of degree $\leq p$, and
\begin{equation}
  E_p(w) = (1-w)\exp\!\bigl(w + w^2/2 + \cdots + w^p/p\bigr)
  \label{eq:weierstrass-primary}
\end{equation}
is the Weierstrass primary factor of genus~$p$.
\end{theorem}

\begin{corollary}[Order-1 factorization]
For entire functions of order~$\rho = 1$ (genus $p = 1$):
\begin{equation}
  f(z) = z^m\, e^{az+b}\, \prod_{n=1}^\infty
  \left(1 - \frac{z}{z_n}\right)\,
  e^{z/z_n}\,,
  \label{eq:hadamard-order1}
\end{equation}
where $a,b\in\mathbb{C}$.
This is the form relevant for $F_1(z)$ and $F_2(z)$.
\end{corollary}


\subsection{Growth from Taylor coefficients}

\begin{proposition}[Order from coefficients]
\label{prop:order-coefficients}
\exactlit\;
\fromlit{Boas, Thm.~2.2.2}
If $f(z) = \sum a_n z^n$ is entire, then
\begin{equation}
  \rho = \limsup_{n\to\infty}
  \frac{n\log n}{-\log|a_n|}\,.
  \label{eq:order-from-coeffs}
\end{equation}
\end{proposition}

\begin{remark}[Application to $\phi$]
For $\phi(z)$, the $n$-th coefficient is
$a_n = (-1)^n n!/(2n+1)!$.  By Stirling,
$|a_n| \sim \sqrt{\pi/(2n)}\cdot 4^{-n}/\sqrt{n}$, so
$-\log|a_n| \sim n\log 4 + O(\log n)$ and
\begin{equation}
  \rho_\phi = \lim_{n\to\infty}\frac{n\log n}{n\log 4 + O(\log n)}
  = \lim_{n\to\infty}\frac{\log n}{\log 4} \cdot
    \frac{1}{1 + O(\log n/(n\log 4))} = \infty \cdot 0\,.
\end{equation}
A more careful analysis using the exact formula gives $\rho_\phi = 1$
\fromlit{verified numerically in NT-2 derivation}.
\end{remark}


%%======================================================================
\section{Analyticity of the Master Function}
\label{sec:phi-analyticity}
%%======================================================================

\begin{theorem}[$\phi(z)$ is entire]
\label{thm:phi-entire}
\exactlit\;
The master function $\phi(z)$ defined by \eqref{eq:phi-integral}
is an entire function of order $\rho = 1$ and type $\sigma = 1/4$.
\end{theorem}
\begin{proof}[Proof outline \fromlit{standard complex analysis}]
\textbf{Step 1 (Entirety):}
The integrand $e^{-\xi(1-\xi)z}$ is entire in $z$ for each fixed
$\xi\in[0,1]$, and on any compact $|z|\leq R$ we have the
uniform bound $|e^{-\xi(1-\xi)z}| \leq e^{R/4}$ (since
$\xi(1-\xi)\leq 1/4$).  By Morera's theorem and uniform convergence
of the parameter integral, $\phi(z)$ is entire.

\textbf{Step 2 (Upper bound):}
For $|z| = r$ on the negative real axis ($z = -r$, the worst case):
\begin{equation}
  |\phi(-r)| = \int_0^1 e^{\xi(1-\xi)r}\,\dd\xi
  \leq e^{r/4}\,,
\end{equation}
using $\xi(1-\xi)\leq 1/4$.  Hence $M(r) \leq C\, e^{r/4}$,
giving $\rho \leq 1$ and $\sigma \leq 1/4$.

\textbf{Step 3 (Lower bound):}
On the negative real axis, by Laplace's method
\fromlit{standard asymptotic analysis}:
\begin{equation}
  \phi(-r) = \int_0^1 e^{\xi(1-\xi)r}\,\dd\xi
  \sim \sqrt{\frac{\pi}{r}}\, e^{r/4}\,,
  \qquad r\to +\infty\,.
  \label{eq:phi-negative-asymp}
\end{equation}
The saddle point is at $\xi = 1/2$, giving the exponential factor
$e^{r/4}$ and the algebraic prefactor $\sqrt{\pi/r}$.
Hence $\log M(r) \geq r/4 + O(\log r)$, giving
$\rho \geq 1$ and $\sigma \geq 1/4$.

Combining: $\rho = 1$, $\sigma = 1/4$.
\end{proof}

\verified{Growth order $\rho=1$, type $\sigma=1/4$ confirmed
numerically by sampling $|\phi(Re^{i\theta})|$ at
$R=50,100,200,400,800$ and fitting $\log\log M(R)/\log R$.
All estimates converge to $\rho \approx 1.0$ and
$\sigma \approx 0.25$.}

\begin{proposition}[Positive real axis decay]
\label{prop:phi-pos-real}
For $z = r > 0$ (positive real axis):
\begin{equation}
  \phi(r) = e^{-r/4}\sqrt{\frac{\pi}{r}}\,\erfi\!\left(
    \frac{\sqrt{r}}{2}\right)
  \to 1\quad\text{as }r\to 0^+\,,
  \qquad
  \phi(r) \sim \frac{2}{r}\quad\text{as }r\to+\infty\,.
\end{equation}
Hence $\phi(r) > 0$ for all $r \geq 0$ (monotonically decreasing
from $\phi(0)=1$ toward $0^+$).
\end{proposition}

\begin{remark}[Physical significance]
The positive real axis corresponds to spacelike momenta
$k^2 > 0$ in Euclidean signature.  The decay of $\phi(r)$
ensures that the form factors approach their local
(Seeley--DeWitt) limits at high momenta, providing the UV
modification mechanism of the spectral action.
\end{remark}


%%======================================================================
\section{Per-Spin Form Factors and Apparent Poles}
\label{sec:form-factors}
%%======================================================================

\subsection{CZ form factor expressions}

All per-spin form factors from the NT-1 derivation are rational
combinations of $\phi(z)$ and $1/z$ \fromlit{NT-1, Theorem~4.1;
CZ Eqs.~(16)--(21)}.  We collect them here for the pole
cancellation analysis.

\begin{convention}[Form factor expressions]
\label{conv:form-factors}
\oursign\;
Using the conventions from NT-1:
\begin{align}
  h_C^{(0)}(z) &= \frac{1}{12z} + \frac{\phi(z)-1}{2z^2}\,,
    \label{eq:hC0}\\[4pt]
  h_R^{(0)}(z;\xi) &= f_{R,\text{bis}}(z) + \xi\,f_{RU}(z)
    + \xi^2\,f_U(z)\,,
    \label{eq:hR0}\\[4pt]
  h_C^{(1/2)}(z) &= \frac{3\phi(z)-1}{6z}
    + \frac{2[\phi(z)-1]}{z^2}\,,
    \label{eq:hC12}\\[4pt]
  h_R^{(1/2)}(z) &= \frac{3\phi(z)+2}{36z}
    + \frac{5[\phi(z)-1]}{6z^2}\,,
    \label{eq:hR12}\\[4pt]
  h_C^{(1)}(z) &= \frac{\phi(z)}{4}
    + \frac{6\phi(z)-5}{6z}
    + \frac{\phi(z)-1}{z^2}\,,
    \label{eq:hC1}\\[4pt]
  h_R^{(1)}(z) &= -\frac{\phi(z)}{48}
    + \frac{11-6\phi(z)}{72z}
    + \frac{5[\phi(z)-1]}{12z^2}\,.
    \label{eq:hR1}
\end{align}
\end{convention}

\begin{remark}[Apparent singularities]
Every form factor has apparent poles at $z=0$ (of order~1 or~2).
These are \textbf{removable singularities}: the functions are
actually regular at $z=0$.  The cancellation mechanism relies on
the identity $\phi(0) = 1$ and the Taylor structure of $\phi(z)$.
\end{remark}


\subsection{Pole cancellation mechanism}

\begin{definition}[Double-zero criterion]
\label{def:double-zero}
A rational expression $f(z) = N(z)/z^k$ is regular at $z=0$
if and only if $N(z) = z^k\, g(z)$ for some analytic $g$.
Equivalently, $N(0) = N'(0) = \cdots = N^{(k-1)}(0) = 0$.
\end{definition}

\begin{proposition}[Key identities for pole cancellation]
\label{prop:key-identities}
The following identities are necessary and sufficient for
regularity of all form factors at $z=0$:
\begin{align}
  \phi(0) &= 1\,,\label{eq:key1}\\
  \phi'(0) &= -\frac{1}{6}\,,\label{eq:key2}\\
  \phi''(0) &= \frac{1}{30}\,.\label{eq:key3}
\end{align}
\end{proposition}
\begin{proof}
Consider $h_C^{(0)}(z) = [z + 6(\phi-1)]/(12z^2)$.
Setting the numerator $N(z) = z + 6[\phi(z)-1]$:
$N(0) = 0 + 6(1-1) = 0$.
$N'(0) = 1 + 6\phi'(0) = 1 + 6(-1/6) = 0$.
Hence $N(z) = z^2\, g(z)$ and $h_C^{(0)}$ is regular.
The other form factors follow analogously (details in
NT-2 derivation document, Propositions~3.1--3.6).
\end{proof}

\verified{All 8 pole cancellation checks (6 per-spin + 2 total SM)
verified at 100-digit precision with $z = 10^{-k}$,
$k = 10, 20, 30, 40, 50$.  Residuals $< 10^{-90}$.}


\subsection{Form factor limiting values}

\begin{proposition}[Zero-momentum limits]
\label{prop:limits}
\fromlit{CZ Table~1; NT-1b Phases 1--3}
\begin{center}
\begin{tabular}{lccc}
\toprule
Form factor & $h(0)$ & Source & Notes \\
\midrule
$h_C^{(0)}(0)$ & $1/120$ & CZ Eq.~(18) & real scalar \\
$h_R^{(0)}(0;\xi)$ & $(\xi-1/6)^2/2$ & CZ Eq.~(19) & $\xi$-dependent \\
$h_C^{(1/2)}(0)$ & $-1/20$ & CZ Eq.~(16) & sign incl.\ in CZ def. \\
$h_R^{(1/2)}(0)$ & $0$ & CZ Eq.~(17) & conformal invariance \\
$h_C^{(1)}(0)$ & $1/10$ & CZ Eq.~(20) & Proca (ghost-corrected) \\
$h_R^{(1)}(0)$ & $0$ & CZ Eq.~(21) & conformal invariance \\
\bottomrule
\end{tabular}
\end{center}
\end{proposition}

\signcorrection{The sign of $h_C^{(1/2)}(0) = -1/20$ (negative) is
correct: the CZ definition for Dirac fields includes a fermionic
sign.  This means $N_D \cdot h_C^{(1/2)}(0) = 22.5 \times (-1/20)
= -135/120$, which correctly produces $\alpha_C = (4 - 135 + 144)/120
= 13/120$ in the SM total.  See NT-1b Phase~3 handoff for the detailed
sign analysis.}


%%======================================================================
\section{Ghost-Free Gravity: Literature Survey}
\label{sec:ghost-free}
%%======================================================================

\subsection{Higher-derivative gravity and the Stelle ghost}

\begin{convention}[Stelle's 1977 result]
\label{conv:stelle}
\fromlit{Stelle, Phys.\ Rev.\ D \textbf{16} (1977) 953}
The most general local quadratic curvature action in 4D is
\begin{equation}
  S = \int\dd^4x\,\sqrt{g}\left[
    \frac{R}{16\pi G} + c_1\, R^2 + c_2\, \Ricsq
  \right].
  \label{eq:stelle-action}
\end{equation}
The graviton propagator in TT gauge has the schematic form
\begin{equation}
  G(k^2) \propto \frac{1}{k^2 + c_2 k^4}\,,
\end{equation}
which decomposes into a massless graviton and a massive spin-2 ghost
(the \textbf{Stelle ghost}) with mass $m_2^2 = 1/c_2$.
The scalar mode mass is $m_0^2 = 1/(3c_1+c_2)$.
\end{convention}

\begin{remark}[The ghost problem]
Stelle showed that the massive spin-2 mode is a ghost
(negative residue).  This makes the local $R + R^2 + R_{\mu\nu}^2$
theory non-unitary, despite its renormalizability.  The nonlocal
structure of the spectral action is designed to avoid this problem.
\end{remark}


\subsection{Nonlocal ghost-free gravity}

\begin{definition}[Ghost-free nonlocal gravity]
\fromlit{Tomboulis hep-th/9702080; BMS hep-th/0508194;
Modesto 1107.2403}
A nonlocal gravitational action of the form
\begin{equation}
  S = \int\dd^4x\,\sqrt{g}\left[
    R + R_{\mu\nu}\,\mathcal{F}_1(\Box)\,R^{\mu\nu}
    + R\,\mathcal{F}_2(\Box)\,R
  \right]
  \label{eq:nonlocal-action}
\end{equation}
is \textbf{ghost-free} if the modified propagator denominators
\begin{equation}
  \Pi_{\text{TT}}(k^2) = 1 + c_2\, k^2\, F_1(k^2/\Lambda^2)\,,
  \qquad
  \Pi_s(k^2) = 1 + (3c_1+c_2)\, k^2\, F_2(k^2/\Lambda^2)
\end{equation}
have \textbf{no real zeros} for $k^2 \geq 0$
\fromlit{BMS Eq.~(7); Modesto Eq.~(4)}.
\end{definition}

\begin{remark}[Sufficient condition: entire form factors]
\fromlit{Tomboulis hep-th/9702080, \S4}
If $\mathcal{F}_i(\Box)$ are chosen such that the propagator
denominators are entire functions with no zeros on the positive
real axis, then the theory has no additional poles (ghosts) beyond
the massless graviton.  Tomboulis showed that choosing
$\mathcal{F}_i(\Box) = e^{H(-\Box/M^2)}$ with $H$ entire and
$H(0) = 0$ ensures this.
\end{remark}

\begin{remark}[BMS infinite-derivative gravity]
\fromlit{BMS hep-th/0508194}
Biswas--Mazumdar--Siegel systematically studied the most general
ghost-free propagator modification of the form
$G^{-1}(k^2) = k^2\,a(-k^2/M^2)$ with $a$ an entire function
having no zeros.  Their key result: the form factor must be
\textbf{exponential of an entire function} to guarantee ghost-freedom.
The simplest choice $a(\Box) = e^{-\Box/M^2}$ was their prototype.
\end{remark}

\begin{remark}[SCT vs.\ engineered models]
\gap\;
The SCT spectral action \textbf{does not} choose the form factors
\emph{a priori} to be ghost-free.  Instead, the form factors
$F_1(z)$, $F_2(z)$ are \emph{derived} from the heat kernel
structure of $\Tr(f(D^2/\Lambda^2))$.  Ghost-freedom is therefore a
\textbf{benchmark to be tested}, not an input.  This is a key
distinction from the BMS/Modesto/Tomboulis approach.

\medskip
The NT-2 task is to verify whether this prediction holds:
are $\Pi_{\text{TT}}(z)$ and $\Pi_s(z)$ zero-free on $z \geq 0$?
\par\smallskip
The repaired NT-4a scan answers this negatively for the current
normalized TT denominator: a positive-real-axis zero is present near
$z \approx 2.41484$.
\end{remark}


%%======================================================================
\section{SM Field Counting Conventions}
\label{sec:sm-counting}
%%======================================================================

\begin{convention}[Standard Model spectrum]
\label{conv:sm-spectrum}
\oursign\;
\fromlit{CPR 0805.2909, Table~1}
\begin{center}
\begin{tabular}{lcc}
\toprule
Field type & Symbol & Value \\
\midrule
Real scalars & $N_s$ & 4 \\
Dirac fermions & $N_D = N_f/2$ & 22.5 \\
Vectors (gauge) & $N_v$ & 12 \\
\bottomrule
\end{tabular}
\end{center}
\end{convention}

\begin{remark}[Half-integer $N_D$]
The value $N_D = 22.5$ arises because $N_f = 45$ counts
\textbf{Weyl} fermion degrees of freedom in the SM (including
3 generations $\times$ 15 Weyl spinors per generation).
The CZ form factors are normalized per Dirac fermion, so we
divide by 2: $N_D = 45/2 = 22.5$.
\end{remark}


%%======================================================================
\section{SM Total Form Factors}
\label{sec:sm-totals}
%%======================================================================

\begin{convention}[Total form factors and local coefficients]
\label{conv:totals}
\oursign\;
\begin{align}
  \alpha_C &= N_s\, h_C^{(0)}(0)
  + N_D\, h_C^{(1/2)}(0)
  + N_v\, h_C^{(1)}(0)
  = \frac{4}{120} - \frac{135}{120} + \frac{144}{120}
  = \frac{13}{120}\,,
  \label{eq:alpha-C}\\[4pt]
  \alpha_R(\xi) &= N_s\, h_R^{(0)}(0;\xi)
  + N_D\, h_R^{(1/2)}(0)
  + N_v\, h_R^{(1)}(0)
  = 4\cdot\frac{(\xi-1/6)^2}{2} + 0 + 0
  = 2\left(\xi-\frac{1}{6}\right)^2.
  \label{eq:alpha-R}
\end{align}
\end{convention}

\begin{remark}[Local limit coefficients]
The local (Stelle-type) coefficients appearing in the propagator
denominators are:
\begin{align}
  c_2 &= 2\alpha_C = \frac{13}{60}\,,\label{eq:c2}\\
  3c_1 + c_2 &= 6\left(\xi - \frac{1}{6}\right)^2.\label{eq:3c1c2}
\end{align}
These are parameter-free ($c_2$) or depend only on the Higgs
non-minimal coupling ($\xi$).
\end{remark}

\verified{$\alpha_C = 13/120$ and $\alpha_R(0) = 1/18$ confirmed
by independent arithmetic and 100-digit mpmath evaluation.
Lean~4 formal proof: \texttt{GhostFree.lean}, Theorem
\texttt{total\_weyl\_beta}.}


%%======================================================================
\section{Propagator Denominator Structure}
\label{sec:propagator}
%%======================================================================

\begin{definition}[Modified propagator denominators]
\label{def:denominators}
\fromlit{BMS Eq.~(7); NT-4a derivation}
The spin-2 (TT) and spin-0 (scalar) propagator denominators are:
\begin{align}
  \Pi_{\text{TT}}(z) &= 1 + \frac{13}{60}\, z\, \hat{F}_1(z)\,,
    \label{eq:Pi-TT}\\
  \Pi_s(z;\xi) &= 1 + 6\left(\xi-\frac{1}{6}\right)^2
    z\, \hat{F}_2(z;\xi)\,,
    \label{eq:Pi-s}
\end{align}
where $\hat{F}_i(z) = F_i(z)/F_i(0)$ is the normalized (unit at
origin) shape function.
\end{definition}

\begin{remark}[Ghost-freedom criterion for SCT]
If $F_1(z)$ and $F_2(z)$ are entire, and $\Pi_{\text{TT}}(z)$,
$\Pi_s(z)$ have no real positive zeros, then the SCT graviton
propagator is ghost-free.  This is the central benchmark imported from
the nonlocal-gravity literature.
\par\smallskip
In the repaired project state this benchmark is \emph{not} satisfied by
the present normalized TT denominator, which crosses zero on the
positive real axis near $z \approx 2.41484$.
\end{remark}

\begin{remark}[Scalar mode decoupling at $\xi = 1/6$]
At conformal coupling $\xi = 1/6$: $\alpha_R = 0$ and
$\Pi_s = 1$ identically.  The scalar mode has the same propagator
as in GR (no modification), and decouples from the ghost-freedom
analysis entirely.  Only $\Pi_{\text{TT}}$ needs to be checked.
\end{remark}


%%======================================================================
\section{Growth Estimates: Extraction Protocol}
\label{sec:growth}
%%======================================================================

\begin{convention}[Numerical growth estimation]
\label{conv:growth}
\oursign\;
To numerically verify the order and type of $F_1(z)$ and $F_2(z)$,
we use the following protocol:
\begin{enumerate}[nosep]
  \item Sample $|F_i(Re^{i\theta})|$ at radii
    $R \in \{50,100,200,400,800\}$ and five angles
    uniformly distributed in $[0,\pi]$.
  \item Compute the order proxy at each radius:
    \[
      \rho_{\mathrm{proxy}} \;=\; \frac{\log\log M(R)}{\log R}\,.
    \]
  \item Compute the type proxy:
    \[
      \sigma_{\mathrm{proxy}} \;=\; \frac{\log M(R)}{R^{\rho}}\,.
    \]
  \item Verify convergence to $\rho = 1$, $\sigma = 1/4$.
\end{enumerate}
\end{convention}

\begin{remark}[Why $\sigma = 1/4$]
The exponential type $\sigma = 1/4$ of $\phi(z)$ comes from the
maximum of $\xi(1-\xi)$ on $[0,1]$, which is $1/4$ at $\xi = 1/2$.
Since $F_1$ and $F_2$ are finite linear combinations of $\phi$
divided by powers of $z$, they inherit the same growth:
$|F_i(z)| \leq C |z|^{-k} e^{|z|/4}$ for suitable $k$.
\end{remark}


%%======================================================================
\section{Zero Structure: What the Literature Predicts}
\label{sec:zeros}
%%======================================================================

\begin{remark}[Positive real axis]
\fromlit{BMS hep-th/0508194, \S3}
For ghost-freedom, we need $\Pi_{\text{TT}}(z) \neq 0$ and
$\Pi_s(z) \neq 0$ for all $z \geq 0$.  At $z = 0$ both equal~1.
This is the \emph{criterion}; it is not automatic from the sign of the
master function alone, because the full denominator depends on the
complete normalized form factor, not merely on the positivity of
$\phi(r)$.
\par\smallskip
The repaired NT-4a scan shows that the current normalized
$\Pi_{\text{TT}}(z)$ does develop a positive-real-axis zero near
$z \approx 2.41484$. Therefore the literature ghost-freedom criterion is
not satisfied in the present normalization.
\end{remark}

\begin{remark}[Complex zeros]
\gap\;
The zeros of $\Pi_{\text{TT}}(z)$ off the real axis are
\textbf{not} directly relevant for Euclidean ghost-freedom,
but they will matter for the Lorentzian continuation (MR-1 task).
These zeros should be tracked and catalogued in the NT-2 analysis
for future use.

\medskip
The literature on \emph{infinite-derivative gravity}
\fromlit{Tomboulis; BMS; Modesto} discusses conditions under which
the entire function $a(\Box)$ can be chosen to avoid complex zeros
entirely (by taking $a = e^{H}$ with $H$ entire and real on the
real axis).  For the SCT spectral action, the form factors are
\emph{not} of this simple exponential form, so complex zeros are
expected to exist.
\end{remark}


%%======================================================================
\section{Conventions Summary}
\label{sec:summary}
%%======================================================================

\begin{center}
\begin{tabular}{lll}
\toprule
Quantity & Convention & Status \\
\midrule
Master function $\phi(z)$ & $\int_0^1 e^{-\xi(1-\xi)z}\,\dd\xi$
  & NT-1 inherited \\
Form factors $h_{C,R}^{(s)}$ & CZ notation with $E = -R/4$ (Dirac)
  & NT-1 inherited \\
SM counting & $N_s=4$, $N_D=22.5$, $N_v=12$
  & CPR 0805.2909 \\
Entire function order/type & $\rho,\sigma$ via Boas definitions
  & Standard \\
Hadamard factorization & Genus $p = \lfloor\rho\rfloor$
  & Standard \\
Ghost-freedom & $\Pi(z) \neq 0$ for $z \geq 0$
  & BMS criterion \\
Propagator denominators & $\Pi = 1 + c\cdot z\cdot\hat{F}(z)$
  & NT-4a convention \\
\bottomrule
\end{tabular}
\end{center}


%%======================================================================
\section{Gap Analysis}
\label{sec:gaps}
%%======================================================================

\begin{enumerate}[label=\textbf{G\arabic*.}]
\item \textbf{Lorentzian continuation:}
  All results in NT-2 are strictly Euclidean.  The physical
  ghost-freedom criterion requires Wick rotation to Lorentzian
  signature, where the propagator poles move to $k^2 = -m^2$
  (timelike momenta).  This is deferred to MR-1.

\item \textbf{Beyond one loop:}
  The form factors $F_1$, $F_2$ are one-loop quantities.
  Higher-loop corrections are $O(G/\Lambda^2)$ and are beyond
  the scope of NT-2.

\item \textbf{Beyond $O(\mathcal{R}^2)$:}
  The entire-function analysis applies to the
  curvature-squared truncation.  Cubic and higher terms in the
  heat kernel expansion are not included.

\item \textbf{Real and complex zero structure:}
  The present normalized TT denominator already develops a
  positive-real-axis zero, so positive-real-axis ghost-freedom is
  \emph{not} established. Beyond that real zero, the full complex zero
  structure of $\Pi(z)$ is only partially explored numerically. A
  rigorous bound on the number and location of complex zeros would
  require deeper analysis (Hadamard genus, Jensen's formula, etc.).
\end{enumerate}


%%======================================================================
\begin{thebibliography}{99}
%%======================================================================

\bibitem{CZ2012}
  A.~Codello, A.~Zanusso,
  ``On the non-local heat kernel expansion,''
  J.\ Math.\ Phys.\ \textbf{54} (2013) 013513
  [arXiv:1203.2034 [hep-th]].

\bibitem{CPR2009}
  A.~Codello, R.~Percacci, C.~Rahmede,
  ``Investigating the ultraviolet properties of gravity with a
  Wilsonian renormalization group equation,''
  Annals Phys.\ \textbf{324} (2009) 414
  [arXiv:0805.2909 [hep-th]].

\bibitem{BV1987}
  A.~O.~Barvinsky, G.~A.~Vilkovisky,
  ``Beyond the Schwinger-DeWitt technique: Converting loops into
  trees and in-in currents,''
  Nucl.\ Phys.\ B \textbf{282} (1987) 163.

\bibitem{BV1990a}
  A.~O.~Barvinsky, G.~A.~Vilkovisky,
  ``Covariant perturbation theory (II): Second order in the
  curvature,''
  Nucl.\ Phys.\ B \textbf{333} (1990) 471.

\bibitem{Vassilevich2003}
  D.~V.~Vassilevich,
  ``Heat kernel expansion: user's manual,''
  Phys.\ Rept.\ \textbf{388} (2003) 279
  [hep-th/0306138].

\bibitem{Avramidi2000}
  I.~G.~Avramidi,
  \emph{Heat Kernel and Quantum Gravity},
  Lecture Notes in Physics \textbf{M64},
  Springer, 2000.

\bibitem{Tomboulis1997}
  E.~T.~Tomboulis,
  ``Super-renormalizable gauge and gravitational theories,''
  hep-th/9702080 (1997).

\bibitem{BMS2006}
  T.~Biswas, A.~Mazumdar, W.~Siegel,
  ``Bouncing universes in string-inspired gravity,''
  JCAP \textbf{0603} (2006) 009
  [hep-th/0508194].

\bibitem{Modesto2012}
  L.~Modesto,
  ``Super-renormalizable quantum gravity,''
  Phys.\ Rev.\ D \textbf{86} (2012) 044005
  [arXiv:1107.2403 [hep-th]].

\bibitem{Stelle1977}
  K.~S.~Stelle,
  ``Renormalization of higher-derivative quantum gravity,''
  Phys.\ Rev.\ D \textbf{16} (1977) 953.

\bibitem{Hadamard1893}
  J.~Hadamard,
  ``\'{E}tude sur les propri\'{e}t\'{e}s des fonctions enti\`{e}res
  et en particulier d'une fonction consid\'{e}r\'{e}e par Riemann,''
  J.\ Math.\ Pures Appl.\ (4) \textbf{9} (1893) 171--215.

\bibitem{Boas1954}
  R.~P.~Boas,
  \emph{Entire Functions},
  Academic Press, New York, 1954.

\bibitem{CC1997}
  A.~H.~Chamseddine, A.~Connes,
  ``The spectral action principle,''
  Commun.\ Math.\ Phys.\ \textbf{186} (1997) 731
  [hep-th/9606001].

\end{thebibliography}

\end{document}
